\documentclass[11pt,a4paper,oneside]{report}
\usepackage[utf8]{inputenc}
\usepackage[T1]{fontenc}
\usepackage[brazil]{babel}
\usepackage[a4paper,margin=2.5cm]{geometry}
\usepackage{lmodern}
\usepackage{microtype}
\usepackage{array,booktabs,longtable,tabularx}
\usepackage{amsmath}
\usepackage{xcolor}
\usepackage{enumitem}
\usepackage{hyperref}
\definecolor{accent}{HTML}{17476B}
\hypersetup{colorlinks=true,linkcolor=accent,urlcolor=accent,
  pdftitle={Projeto Finanças: aplicação, home server e rede doméstica},
  pdfsubject={Relatório técnico e financeiro em desenvolvimento}}
\setlength{\parindent}{0pt}
\setlength{\parskip}{6pt}
\setlength{\emergencystretch}{3em}
\renewcommand{\arraystretch}{1.2}
\setlist{nosep,leftmargin=*}
\newcolumntype{Y}{>{\raggedright\arraybackslash}X}
\newcommand{\pendente}{\textit{A definir}}
\newcommand{\naoinformado}{\textit{Não informado}}
\newcommand{\dinheiro}[1]{R\$\,#1}
% Atualizar estes campos e manter as revisões futuras neste mesmo arquivo.
\newcommand{\versaorelatorio}{0.5}
\newcommand{\datarelatorio}{8 de outubro de 2026}
\newcommand{\autorrelatorio}{Autor: a definir}

\begin{document}
\begin{titlepage}
\centering
\vspace*{2cm}
{\Large RELATÓRIO DE DESENVOLVIMENTO\par}
\vspace{1.2cm}
{\Huge\bfseries\color{accent} Projeto Finanças\par}
\vspace{0.7cm}
{\Large Aplicação financeira, home server\\e rede doméstica\par}
\vspace{1.5cm}
{\large Documentação técnica e levantamento financeiro\par}
\vfill
{\large \autorrelatorio\par}
\vspace{0.5cm}
Versão \versaorelatorio{} --- documento em desenvolvimento\par
\datarelatorio\par
\end{titlepage}

\chapter*{Resumo}
\addcontentsline{toc}{chapter}{Resumo}
Este relatório acompanha a evolução do Projeto Finanças, identificado como
FinanceHub na documentação da aplicação, e sua futura hospedagem em um servidor
residencial. O escopo reúne o desenvolvimento do software, o reaproveitamento e a
aquisição de componentes, a montagem do computador, a instalação do Proxmox VE,
a implantação dos serviços
e a organização da rede doméstica.

Na data desta versão, a aplicação está em fase de finalização e a aquisição de
componentes para o servidor está começando, conforme informado pelo responsável.
A instalação definitiva, a configuração da rede e os resultados dos testes no
servidor ainda precisam ser registrados. O levantamento financeiro separa o
patrimônio já disponível, os novos desembolsos e as despesas recorrentes, sem
atribuir valores a compras ou medições que não foram informadas.
O orçamento inicial informado é de \dinheiro{15.000,00}. A primeira lista de
componentes para o computador principal soma \dinheiro{11.100,00}.
Com um rack de piso de 12U e dois gabinetes rack 4U, estimados em
\dinheiro{1.744,00}, e um switch de \dinheiro{400,00}, o total dos itens com
preço informado passa a \dinheiro{13.244,00}, com margem preliminar de
\dinheiro{1.756,00} antes de fretes e itens adicionais. Está previsto cabeamento
de rede para dez máquinas, com metragens e custo ainda a levantar.
As peças antigas serão reaproveitadas no servidor. A correção da memória DDR4
indicada na plataforma AM5 ainda está pendente.
Os valores são de planejamento, não gastos comprovados.

\textbf{Palavras-chave:} finanças pessoais; servidor doméstico; Proxmox VE; Docker;
rede de computadores; orçamento; documentação.

\tableofcontents
\clearpage

\chapter{Introdução e objetivos}
\section{Contexto}
O projeto parte de uma aplicação de gerenciamento financeiro e avança para a
construção de uma infraestrutura doméstica capaz de hospedá-la e de executar
outros serviços. O planejamento do home server inclui Proxmox VE como plataforma
de virtualização, Docker, automações e serviços domésticos.
Os demais serviços e o volume de utilização serão
definidos durante a implantação.

A infraestrutura e a aplicação fazem parte de um mesmo processo: desenvolver
o sistema, preparar o equipamento, disponibilizar o acesso e estabelecer
rotinas de operação, manutenção e recuperação dos dados.

\section{Objetivo geral}
Documentar a construção e a operação do Projeto Finanças em um home server,
integrando decisões de software, hardware, rede e viabilidade financeira.

\section{Objetivos específicos}
\begin{enumerate}
\item Registrar a arquitetura e a finalização da aplicação financeira.
\item Levantar os componentes existentes e verificar sua compatibilidade.
\item Documentar as compras e a montagem física do servidor.
\item Documentar a instalação do Proxmox VE e a configuração do ambiente virtualizado.
\item Registrar a escolha do sistema das máquinas virtuais e a implantação dos serviços.
\item Mapear e configurar a rede doméstica e o acesso autorizado à aplicação.
\item Acompanhar investimento inicial, despesas recorrentes e desvios do orçamento.
\item Reunir evidências de testes, manutenção e restauração dos backups.
\end{enumerate}

\section{Limites desta versão}
Esta versão constitui a base do relatório. Não representa uma instalação
concluída nem comprova funcionamento contínuo, desempenho, consumo elétrico ou
segurança do servidor real. A lista final de compras depende da confirmação
da compatibilidade e dos itens adicionais para o computador principal,
o servidor e a rede.

\chapter{Método de documentação e andamento}
\section{Origem das informações}
As informações iniciais foram obtidas dos relatos do responsável e da leitura
dos arquivos \path{README.md}, \path{backend/README.md},
\path{frontend/ARCHITECTURE.md} e \path{deploy/README.md} do repositório.
Especificações externas de hardware são identificadas nas referências.
Informações de documentação descrevem o que o projeto declara ou prepara;
resultados operacionais exigem registros da execução no equipamento de destino.

\section{Classificação dos registros}
\begin{description}
\item[Informado:] dado fornecido pelo responsável, ainda sujeito a conferência física.
\item[Documentado:] recurso ou procedimento descrito no repositório.
\item[Planejado:] atividade ou arquitetura proposta, ainda sem execução confirmada.
\item[Validado:] resultado acompanhado de data, método e evidência verificável.
\item[Pendente:] informação ausente ou decisão ainda aberta.
\end{description}

\section{Situação em \datarelatorio}
\begin{tabularx}{\textwidth}{@{}p{4.1cm}YY@{}}
\toprule
Frente & Situação & Próximo registro \\
\midrule
Projeto Finanças & Em finalização, segundo o responsável & Versão adotada e resultados da validação final \\
Componentes & Computador principal novo; peças antigas para o servidor & Compatibilidade e modelos completos \\
Montagem física & Sem conclusão registrada & Sequência de montagem e testes \\
Proxmox VE & Instalação planejada & Versão, armazenamento, rede e máquina virtual \\
Home server & Implantação planejada & Sistema da máquina virtual e serviços instalados \\
Rede doméstica & Switch definido e cabeamento para dez máquinas previsto & Percursos, metragens, preços e endereçamento \\
Controle financeiro & Itens com preço: \dinheiro{13.244,00} para orçamento de \dinheiro{15.000,00} & Cabeamento, fretes, demais itens e despesas realizadas \\
\bottomrule
\end{tabularx}

\section{Regra de atualização}
Cada etapa deverá registrar a data, o objetivo, a decisão, sua justificativa,
os itens envolvidos, o custo, o resultado e a evidência. Fotografias, diagramas,
medições e comprovantes poderão ser incorporados em versões futuras.
Capturas devem ocultar credenciais, chaves, dados bancários e identificadores
pessoais desnecessários ao relato.

\chapter{Desenvolvimento do Projeto Finanças}
\section{Finalidade e estrutura atual}
O FinanceHub é descrito no repositório como um gerenciador financeiro com
backend Flask, frontend React e integração bancária com Pluggy. A documentação
de produção adota SQLite para a persistência dos dados.

\begin{tabularx}{\textwidth}{@{}p{3cm}YY@{}}
\toprule
Camada & Tecnologia documentada & Responsabilidade \\
\midrule
Interface & React e Vite & Login, navegação e apresentação dos dados financeiros \\
API & Flask e SQLAlchemy & Autorização, regras, integração e persistência \\
Dados & SQLite & Armazenamento local da aplicação \\
Integração & Pluggy & Conexão de contas e sincronização bancária \\
Produção & Waitress e Caddy & Execução da API, arquivos estáticos e proxy \\
\bottomrule
\end{tabularx}

No frontend, o fluxo documentado é página, hook, serviço HTTP e API. No backend,
a organização utiliza rotas, serviços, repositórios e modelos, com exceções
descritas no guia técnico. Essa separação auxilia a manutenção e a localização
das responsabilidades do sistema.

\section{Funcionalidades descritas}
\begin{itemize}
\item Autenticação individual e contas criadas pelo administrador.
\item Visualização de contas, transações, saldos e indicadores financeiros.
\item Sincronização manual de transações bancárias por mês.
\item Classificação por categorias e subcategorias e associação de lançamentos.
\item Filtros de período e preferência para ocultar valores na interface.
\end{itemize}
A comprovação da versão final deverá incluir testes desses fluxos com dados de
teste e, quando necessário, validação controlada da integração real.

\section{Proteção dos dados e recuperação}
A documentação descreve sessões JWT em cookies HttpOnly, proteção CSRF e
criptografia AES-256-GCM do conteúdo financeiro de contas e transações.
A criptografia financeira não abrange todas as tabelas e metadados do banco.
O administrador e processos com acesso às chaves podem decifrar os dados.

Os procedimentos de produção preveem separação entre código, configuração,
banco, chaves e backups; cópias periódicas criptografadas; verificação de
integridade; e restauração em diretório separado. Guardar uma cópia no próprio
servidor não cobre a perda desse equipamento: o destino externo e a preservação
das chaves precisam ser definidos e testados.

\section{Finalização e evidências pendentes}
\begin{tabularx}{\textwidth}{@{}p{4cm}YY@{}}
\toprule
Registro & Evidência esperada & Situação nesta versão \\
\midrule
Versão de entrega & Identificador da revisão ou versão & \pendente \\
Validação do backend & Resultado e data dos testes & A registrar para a entrega \\
Validação do frontend & Build, verificações e fluxos no navegador & A registrar para a entrega \\
Integração bancária & Sincronização controlada e resultado & A registrar \\
Backup e recuperação & Arquivo verificado e restauração isolada & A validar no servidor \\
\bottomrule
\end{tabularx}
Os testes históricos descritos no repositório não substituem a validação da
versão efetivamente instalada. Nesta preparação do relatório, não foram
executados testes da aplicação ou alterados seus serviços.

\chapter{Componentes e montagem do computador}
\section{Inventário informado}
O inventário abaixo corresponde às peças existentes informadas inicialmente.
A nova lista de compras é registrada separadamente. O responsável confirmou
o reaproveitamento das peças antigas no servidor, após a atualização do computador
principal; a execução dessa transferência ainda será documentada.
\begin{tabularx}{\textwidth}{@{}p{3cm}p{3.7cm}Y@{}}
\toprule
Componente & Informação disponível & Conferência necessária \\
\midrule
Processador & Intel Core i5-10400F & Confirmar modelo e estado \\
Placa-mãe & B450M & Confirmar fabricante e modelo completo; incompatibilidade aparente \\
Memória & 16 GB DDR4 & Quantidade de módulos, frequência e teste \\
Armazenamento & HD de 1 TB & Modelo, estado e dados já existentes \\
Fonte & 650 W & Fabricante, modelo, idade e estado \\
Gabinete e cooler & \naoinformado & Dimensões, suporte às peças e ventilação \\
Placa de vídeo & \naoinformado & Disponibilidade para instalação e diagnóstico \\
\bottomrule
\end{tabularx}

\section{Pendência de compatibilidade}
O Intel Core i5-10400F usa socket LGA1200 e não possui gráficos integrados,
conforme a documentação Intel \cite{intelcpu,intelf}. A família B450 pertence
à plataforma AMD AM4 \cite{amdb450}. Portanto, o conjunto informado não pode
ser montado como está. O modelo correto deverá ser conferido antes de escolher
uma placa substituta ou comprar outro processador.

A necessidade de manter uma placa de vídeo durante a operação sem monitor
depende do comportamento de inicialização da placa-mãe real e será testada.
O nome B450M sozinho também não identifica todas as especificações da placa.

\section{Nova lista de componentes para compra}
Foi informada uma configuração com AMD Ryzen 7 9700X, water cooler, placa-mãe
ASUS B650M, memória de 32 GB, Radeon RX 9070 XT de 16 GB, SSD de 500 GB,
HD de 2 TB e fonte modular de 850 W, 80 Plus Gold. Os preços de cada item
constam no levantamento financeiro. Essa configuração é destinada ao
\textbf{computador principal}; as peças antigas serão utilizadas no home server.
Não há confirmação de compra, recebimento ou montagem desses itens.

\subsection{Compatibilidade e especificações pendentes}
\textbf{A memória DDR4 de 32 GB (2 módulos de 16 GB), 3200 MHz, da lista não é
compatível com o conjunto Ryzen 7 9700X e B650.} A plataforma utiliza socket
AM5 e memória DDR5 \cite{ryzen9700,asusb650}. O item deve ser corrigido para
um kit DDR5 compatível; modelo, frequência e preço serão confirmados novamente.
O valor de \dinheiro{1.000,00} foi mantido apenas para reproduzir a lista recebida.

ASUS B650M identifica uma família, não o modelo completo. Será necessário
consultar a lista de CPUs suportadas e a versão de BIOS exigida para o 9700X
na placa escolhida. Também permanecem pendentes os modelos do water cooler,
SSD, HD, fonte e placa de vídeo, além do gabinete e das dimensões para montagem.
A Radeon RX 9070 XT está prevista para o computador principal e não faz parte
da configuração planejada do servidor. A placa de vídeo disponível para a
montagem inicial do servidor com o i5-10400F ainda precisa ser identificada.

\section{Rack e gabinetes previstos}
A infraestrutura física passa a incluir um rack de servidor de piso,
padrão de 19 polegadas, 12U, medida de 470 mm informada e estrutura desmontável,
além de dois gabinetes K-mex rack de 19 polegadas, 4U, modelo CR-S350, sem fonte.
Os valores previstos são \dinheiro{944,00} para o rack e \dinheiro{400,00}
por gabinete. A interpretação do preço dos gabinetes como unitário deverá
ser conferida na cotação ou no pedido.

Os dois gabinetes representam 8U de ocupação nominal, restando 4U dos 12U
do rack para outros equipamentos e acessórios. Essa contagem não confirma
o encaixe físico: serão conferidos a profundidade útil do rack, as dimensões
dos gabinetes, o método de fixação e o espaço para cabos e ventilação.
A medida de 470 mm não foi tratada como profundidade útil confirmada.
A destinação de cada gabinete será registrada no plano de montagem.
As fontes necessárias serão contabilizadas conforme reaproveitamento ou compra,
pois os gabinetes informados não incluem fonte.

\section{Itens a avaliar para aquisição}
\begin{itemize}
\item SSD para sistema, serviços e dados ativos: modelo e capacidade a definir.
\item Correção do conjunto processador e placa-mãe, se a incompatibilidade for confirmada.
\item Cooler, ventoinhas, cabos e acessórios de fixação dos gabinetes no rack,
conforme o inventário físico e a conferência dimensional.
\item Meio separado para backups e, conforme orçamento, proteção elétrica.
\item Componentes de rede somente após verificar os equipamentos já existentes.
\end{itemize}
O dimensionamento deverá distinguir os 16 GB do inventário existente e os
32 GB propostos na lista nova, cuja tecnologia precisa ser corrigida. A memória
do servidor parte dos 16 GB DDR4 existentes; os 32 GB da lista nova são para o
computador principal. Uma expansão do servidor será avaliada considerando
Proxmox VE, máquinas virtuais e serviços. Nenhuma aquisição é registrada
como realizada sem confirmação de item, data e valor.

\section{Sequência prevista de montagem}
\begin{enumerate}
\item Conferir modelos, compatibilidade, acessórios e manuais dos componentes.
\item Registrar fotos e o estado inicial das peças, com o equipamento desligado.
\item Instalar processador, cooler e memória conforme os manuais.
\item Fixar placa-mãe, fonte e armazenamento e conectar os cabos.
\item Conferir alimentação, ventilação e conexões antes da primeira partida.
\item Registrar detecção de CPU, memória e discos no firmware.
\item Testar estabilidade, memória, temperaturas e estado dos discos.
\end{enumerate}

\section{Registro da execução}
\begin{tabularx}{\textwidth}{@{}p{4.4cm}Y@{}}
\toprule
Campo & Informação \\
\midrule
Data da montagem & \pendente \\
Configuração efetivamente montada & \pendente \\
Versão do firmware e ajustes & \pendente \\
Problemas e soluções & A registrar durante a montagem \\
Resultados e fotografias & A registrar após a execução \\
\bottomrule
\end{tabularx}

\chapter{Implantação do home server}
\section{Arquitetura prevista}
O planejamento passa a incluir a instalação do Proxmox VE no computador físico
e a criação de uma máquina virtual Linux para executar o Docker e os serviços
domésticos. A versão do Proxmox, a distribuição Linux da máquina virtual e a
alocação de recursos ainda serão definidas. O repositório também contém
procedimentos de produção para Windows com serviços nativos.

O Proxmox VE gerencia máquinas virtuais e contêineres. Para os serviços Docker,
o desenho previsto utiliza uma máquina virtual, conforme a orientação da
documentação oficial para essa forma de implantação \cite{proxmox}.

\begin{center}
\fbox{\parbox{0.88\textwidth}{\centering
Computador físico $\longrightarrow$ Proxmox VE\\[4pt]
Proxmox VE $\longrightarrow$ máquina virtual Linux\\[4pt]
Máquina virtual Linux $\longrightarrow$ Docker e serviços}}
\end{center}

A documentação consultada apresenta como desenho de acesso privado:
\begin{center}
\fbox{\parbox{0.88\textwidth}{\centering
Dispositivo autorizado\\
$\downarrow$ HTTPS privado\\
Tailscale Serve na máquina virtual Linux\\
$\downarrow$ proxy local\\
Caddy: interface React e encaminhamento da API\\
$\downarrow$\\
Waitress / Flask $\longrightarrow$ SQLite e chaves protegidas}}
\end{center}
Esse desenho é uma referência de implantação, ainda sem validação no servidor.
O estado dos arquivos de implantação Docker também deverá ser conferido antes
da instalação: a existência de referências no README não comprova que todos
os artefatos ou testes estejam concluídos.

\section{Instalação e configuração do Proxmox VE}
Esta etapa ocorrerá após a montagem e a validação do computador. Serão
documentados a versão instalada, o disco de destino, os ajustes de virtualização
no firmware, o acesso administrativo, o armazenamento e a configuração da rede.
Os dados existentes nos discos serão conferidos e preservados antes da instalação.

\begin{enumerate}
\item Conferir os requisitos da versão escolhida e o suporte à virtualização.
\item Preparar a mídia de instalação e registrar o disco selecionado.
\item Instalar o Proxmox VE e configurar o endereço de gerenciamento.
\item Preparar o armazenamento e a conexão de rede das máquinas virtuais.
\item Criar a máquina virtual Linux e registrar CPU, memória e espaço alocados.
\item Definir inicialização automática, atualização e rotinas de backup.
\item Testar reinício do host, recuperação da máquina virtual e acesso aos serviços.
\end{enumerate}

\begin{tabularx}{\textwidth}{@{}p{4.4cm}Y@{}}
\toprule
Registro & Situação \\
\midrule
Versão e data de instalação & \pendente \\
Discos e armazenamento do Proxmox & \pendente \\
Endereço de gerenciamento e rede virtual & \pendente \\
Sistema e recursos da máquina virtual & \pendente \\
Backup e restauração da máquina virtual & A validar após a instalação \\
\bottomrule
\end{tabularx}
O backup da máquina virtual será documentado em conjunto com o backup da
aplicação, registrando a finalidade de cada cópia e os testes de recuperação.

\section{Etapas a documentar}
\begin{enumerate}
\item Instalação e configuração do Proxmox VE no computador físico.
\item Criação da máquina virtual Linux e definição dos recursos e da rede.
\item Instalação do sistema da máquina virtual, com versão e justificativa.
\item Preparação de discos, usuários, permissões e atualizações no ambiente virtual.
\item Instalação do Docker na máquina virtual e definição de serviços, redes e volumes persistentes.
\item Preparação ou migração consistente do banco, configuração e chaves.
\item Implantação da aplicação e configuração do acesso privado.
\item Configuração dos backups, destino externo e monitoramento.
\item Testes de acesso, reinício, isolamento dos usuários e recuperação.
\end{enumerate}

\section{Separação do estado persistente}
O relatório deverá registrar onde ficam o banco, os arquivos de configuração,
as chaves, os backups e os logs. Esses caminhos serão definidos na instalação.
Credenciais não deverão ser incorporadas às imagens nem aos arquivos públicos
do frontend. O procedimento de migração precisa preservar as chaves existentes
quando houver dados criptografados.

\section{Critérios propostos para entrada em operação}
\begin{itemize}
\item Aplicação acessível pelos dispositivos autorizados e login funcional.
\item Dados e permissões preservados após reiniciar o servidor.
\item Serviços iniciados de forma consistente após a reinicialização.
\item Máquina virtual e serviços recuperados após reiniciar o host Proxmox VE.
\item Backup recente verificado e restauração concluída em ambiente separado.
\item Consumo de recursos e falhas observáveis por logs e medições.
\end{itemize}
Datas, responsáveis e resultados desses critérios serão preenchidos após os testes.

\chapter{Rede doméstica}
\section{Levantamento inicial}
A rede existente ainda precisa ser mapeada. Foi incluído no planejamento um
switch TP-Link TL-SG1016D por \dinheiro{400,00} e cabeamento para dez máquinas.
Serão levantados modem ou ONT, roteador, pontos de acesso, percursos dos cabos
e dispositivos que acessarão o servidor.

\begin{tabularx}{\textwidth}{@{}p{4cm}Y@{}}
\toprule
Informação & Situação \\
\midrule
Provedor e plano de internet & \naoinformado \\
Modem/ONT e roteador & Modelos e funções a confirmar \\
Switch & TP-Link TL-SG1016D previsto; instalação pendente \\
Pontos de acesso Wi-Fi & Existência e modelos a confirmar \\
Cabeamento e pontos disponíveis & Dez conexões previstas; percursos, metragens e custo pendentes \\
Local do servidor & Alimentação, ventilação e conexão a definir \\
Sub-rede e DHCP & Configuração atual a levantar \\
Dispositivos e usuários autorizados & Lista a elaborar \\
\bottomrule
\end{tabularx}

\section{Switch e cabeamento para dez máquinas}
O TL-SG1016D é um switch não gerenciável com 16 portas RJ45 de
10/100/1000 Mbps, conforme a documentação do fabricante \cite{tplinkswitch}.
Sua função prevista é conectar as máquinas cabeadas à rede doméstica.
Serão registrados a versão de hardware adquirida, o local de instalação,
a identificação dos cabos e a ocupação de cada porta.

Dez máquinas e uma conexão ao roteador ocupam onze portas, deixando cinco
disponíveis nessa configuração. Se o servidor não estiver entre as dez máquinas,
sua conexão ocupará uma porta adicional. O inventário de portas será ajustado
quando a relação dos equipamentos for definida.

O cabeamento é planejado inicialmente com dez conexões, uma por máquina,
usando categoria 6 como proposta a confirmar no levantamento. Para cada ponto,
serão medidos o percurso até o switch, a folga necessária e o comprimento final.
O total de metros será a soma dos dez percursos; nenhuma metragem foi presumida.

\begin{tabularx}{\textwidth}{@{}p{4.5cm}Y@{}}
\toprule
Item de cabeamento & Dimensionamento previsto \\
\midrule
Conexões às máquinas & Dez pontos, identificados de P01 a P10 \\
Cabos de rede & Dez percursos; categoria e metragens a confirmar \\
Ligação switch--roteador & Uma conexão adicional; verificar cabo disponível \\
Terminações e cabos de conexão & Conforme instalação direta ou pontos com tomadas \\
Organização e fixação & Etiquetas, organizadores e canaletas conforme os percursos \\
Preço e instalação & A cotar após definir comprimentos e forma de montagem \\
\bottomrule
\end{tabularx}
Tomadas, patch panel e cabos de conexão serão dimensionados se for adotado
cabeamento estruturado. O levantamento financeiro mantém o cabeamento como
custo pendente, sem tratá-lo como gratuito. Cada conexão será identificada
e testada antes da entrada em operação.

\section{Topologia conceitual}
\begin{center}
\fbox{\parbox{0.88\textwidth}{\centering
Internet $\longrightarrow$ modem/ONT $\longrightarrow$ roteador\\[4pt]
Roteador $\longrightarrow$ switch TP-Link TL-SG1016D\\[4pt]
Switch $\longrightarrow$ dez máquinas por cabo\\[4pt]
Servidor: entre as dez máquinas ou em porta adicional, a confirmar\\[4pt]
Roteador $\longrightarrow$ dispositivos Wi-Fi}}
\end{center}
O desenho é conceitual: modem e roteador podem ser um único equipamento.
O diagrama real será atualizado com modelos, portas, conexões e locais.

\section{Plano de configuração}
Serão registrados os endereços do host Proxmox VE e da máquina virtual,
a conexão entre a rede física e a virtual, a faixa DHCP, os nomes utilizados,
as regras de acesso e o método de administração remota. Uma reserva DHCP para
o servidor poderá ser definida conforme a rede existente. Segmentação, rede
de visitantes e isolamento de dispositivos de automação serão avaliados de
acordo com os recursos do roteador e as necessidades identificadas.
O switch escolhido não oferece gerenciamento de VLANs; qualquer segmentação
deverá considerar a capacidade dos demais equipamentos e o desenho final da rede.

Para a aplicação financeira, a documentação atual prevê acesso privado por
Tailscale Serve. Qualquer mudança desse desenho exigirá registrar a justificativa
e os testes correspondentes. Endereços, nomes e permissões reais estão pendentes.

\section{Testes previstos}
\begin{itemize}
\item Conectividade do servidor e manutenção do endereço após reinício.
\item Resolução de nomes e acesso autorizado à aplicação.
\item Comportamento do acesso remoto e das restrições aplicadas.
\item Taxa de transferência local, método de medição e equipamentos envolvidos.
\item Identificação e teste das dez conexões e da ligação ao roteador.
\item Recuperação da conexão após reiniciar roteador e servidor.
\end{itemize}

\chapter{Levantamento e acompanhamento financeiro}
\section{Critérios do levantamento}
O controle financeiro terá três grupos: peças já disponíveis, investimento
novo e despesas recorrentes. O preço histórico de uma peça existente será
registrado apenas quando conhecido; ele não será contado novamente como compra
do home server. Valor de mercado ou de reposição, se levantado, ficará separado.

Uma peça reaproveitada pode ter custo incremental de aquisição igual a zero,
mas somente após confirmar seu uso e eventuais reparos. Campos sem informação
não significam custo zero. O responsável informou um orçamento inicial de
\textbf{\dinheiro{15.000,00}}. Os componentes do computador principal somam
\dinheiro{11.100,00}; rack e dois gabinetes somam \dinheiro{1.744,00};
o switch acrescenta \dinheiro{400,00}, resultando em
\textbf{\dinheiro{13.244,00}} de itens com preço informado.
O cabeamento para dez máquinas foi incluído com custo a cotar.
Os preços são fornecidos pelo responsável,
sem validação independente de mercado, fornecedor ou disponibilidade.
O total pago permanece indeterminado, pois não foram informados comprovantes
ou valores efetivamente desembolsados.

O limite é tratado como orçamento global de investimento. A lista de
\dinheiro{11.100,00} pertence ao computador principal; o servidor utilizará
as peças antigas. Rack e gabinetes foram registrados como infraestrutura física.
Demais adequações do servidor, rede, backups e reserva serão dimensionados
dentro da margem preliminar de \dinheiro{1.756,00}, que também precisará cobrir
o cabeamento ainda sem preço.
A estimativa é provisória, sobretudo pela correção necessária da memória DDR4.

\section{Limites e indicadores}
\begin{tabularx}{\textwidth}{@{}p{6.2cm}Y@{}}
\toprule
Indicador & Situação \\
\midrule
Orçamento global informado & \textbf{\dinheiro{15.000,00}} \\
Subtotal previsto do computador principal & \textbf{\dinheiro{11.100,00}} \\
Subtotal previsto de rack e gabinetes & \textbf{\dinheiro{1.744,00}} \\
Switch previsto & \dinheiro{400,00} \\
Cabeamento para dez máquinas & A cotar; não incluído no total abaixo \\
Total previsto dos itens com preço & \textbf{\dinheiro{13.244,00}} \\
Parcela do orçamento comprometida na estimativa & 88,29\% \\
Margem preliminar após os itens com preço & \textbf{\dinheiro{1.756,00}} (11,71\%) \\
Novas despesas do servidor & Reaproveitamento previsto; adequações a levantar \\
Parcela destinada à rede e aos backups & A distribuir dentro do orçamento global \\
Reserva para imprevistos & Percentual ou valor a definir \\
Total pago confirmado & Indeterminado; pagamentos não informados \\
Total previsto para concluir & Total listado mais correções e itens ainda não levantados \\
Desembolso restante & Depende de compras e compromissos registrados \\
Custo mensal e anual de operação & Depende de medição e tarifas \\
\bottomrule
\end{tabularx}

\section{Patrimônio disponível}
\begin{tabularx}{\textwidth}{@{}p{4.3cm}p{3cm}Y@{}}
\toprule
Item informado & Valor histórico & Tratamento no projeto \\
\midrule
i5-10400F & \naoinformado & Reaproveitamento sujeito à confirmação \\
Placa B450M & \naoinformado & Modelo e compatibilidade pendentes \\
Memória DDR4, 16 GB & \naoinformado & Reaproveitamento a validar \\
HD, 1 TB & \naoinformado & Reaproveitamento a validar \\
Fonte, 650 W & \naoinformado & Reaproveitamento a validar \\
\bottomrule
\end{tabularx}

\section{Computador principal: componentes e preços informados}
\begin{longtable}{@{}p{0.7cm}p{3.5cm}p{5cm}r@{}}
\toprule
ID & Componente & Especificação informada & Valor \\
\midrule
\endfirsthead
\toprule
ID & Componente & Especificação informada & Valor \\
\midrule
\endhead
A01 & Processador & AMD Ryzen 7 9700X & \dinheiro{1.700,00} \\
A02 & Water cooler & Modelo e dimensões não informados & \dinheiro{200,00} \\
A03 & Placa-mãe & ASUS B650M; modelo completo a confirmar & \dinheiro{750,00} \\
A04 & Memória & DDR4, 32 GB (2 x 16 GB), 3200 MHz; incompatível com a plataforma & \dinheiro{1.000,00} \\
A05 & Placa de vídeo & AMD Radeon RX 9070 XT, 16 GB & \dinheiro{5.600,00} \\
A06 & SSD & 500 GB; modelo e interface a confirmar & \dinheiro{600,00} \\
A07 & HD & 2 TB; modelo a confirmar & \dinheiro{750,00} \\
A08 & Fonte & 850 W, 80 Plus Gold, modular; modelo a confirmar & \dinheiro{500,00} \\
\midrule
\multicolumn{3}{@{}l}{\textbf{Subtotal informado}} & \textbf{\dinheiro{11.100,00}} \\
\bottomrule
\end{longtable}
O preço da memória foi considerado como valor do kit completo, não por módulo.
Os demais preços foram tratados como o valor de uma unidade de cada item.
Esta tabela reproduz o planejamento recebido, sem registrar compras realizadas.
Fretes, encargos e descontos não foram discriminados e não se presume que
estejam incluídos. Cada aquisição deverá identificar fornecedor, data e garantia.

\section{Infraestrutura física: rack e gabinetes}
\begin{tabularx}{\textwidth}{@{}Ycrr@{}}
\toprule
Item & Qtd. & Valor unitário & Subtotal \\
\midrule
Rack servidor de piso, 19 polegadas, 12U, 470 mm, desmontável
& 1 & \dinheiro{944,00} & \dinheiro{944,00} \\
Gabinete servidor K-mex rack 19 polegadas, 4U, sem fonte, CR-S350
& 2 & \dinheiro{400,00} & \dinheiro{800,00} \\
\midrule
\multicolumn{3}{@{}l}{\textbf{Subtotal da infraestrutura física}}
& \textbf{\dinheiro{1.744,00}} \\
\bottomrule
\end{tabularx}
O valor de \dinheiro{400,00} foi considerado por gabinete; o custo das duas
unidades foi calculado como \dinheiro{800,00}. Essa premissa permanece explícita
até a conferência do preço. Não foram informados frete, fornecedor, garantia
ou confirmação de pagamento desses itens.

\section{Rede: switch e cabeamento}
\begin{tabularx}{\textwidth}{@{}Ycrr@{}}
\toprule
Item & Qtd. & Valor unitário & Subtotal \\
\midrule
Switch TP-Link TL-SG1016D, 16 portas, 10/100/1000 Mbps
& 1 & \dinheiro{400,00} & \dinheiro{400,00} \\
Cabeamento de rede para máquinas, com metragens a levantar
& 10 pontos & \pendente & \pendente \\
\midrule
\multicolumn{3}{@{}l}{\textbf{Subtotal da rede com preço informado}}
& \textbf{\dinheiro{400,00}} \\
\bottomrule
\end{tabularx}
O subtotal inclui somente o switch. Cabos, terminações, acessórios e instalação
permanecem pendentes de cotação. O cabo da ligação ao roteador será incluído
se não houver um disponível para reaproveitamento. Não há confirmação de compra
ou pagamento desses itens.

\section{Análise do orçamento e itens adicionais}
\begin{align*}
I_{\mathrm{PC}} &= \dinheiro{11.100,00},\\
I_{\mathrm{infra}} &= 944 + 2 \times 400 = \dinheiro{1.744,00},\\
I_{\mathrm{switch}} &= \dinheiro{400,00},\\
I_{\mathrm{lista}} &= 11100 + 1744 + 400 = \dinheiro{13.244,00},\\
S_{\mathrm{preliminar}} &= \dinheiro{15.000,00} - \dinheiro{13.244,00}
= \dinheiro{1.756,00}.
\end{align*}
A margem preliminar não é um saldo de caixa ou confirmação de economia:
precisa absorver o cabeamento de dez máquinas, a diferença de preço da memória compatível, fretes,
acessórios e demais itens ainda ausentes do levantamento. A GPU representa
aproximadamente 50,5\% do subtotal do computador principal. Esse valor não
constitui gasto com placa de vídeo do servidor.

Ainda serão levantados ventoinhas, fixações do rack, cabos, armazenamento de backup
separado, demais equipamentos de rede e proteção elétrica, conforme a necessidade e
o inventário existente. Não há valores atribuídos a esses itens. Nenhum custo
de uma configuração nova será somado novamente como peça reaproveitada.

\section{Registro de cotações e compras}
Para cada item, usar o custo completo:
\begin{equation}
C_{\mathrm{item}} = q p + f + a - d,
\end{equation}
em que $q$ é a quantidade, $p$ o preço unitário, $f$ o frete, $a$ os encargos
adicionais não incluídos no preço e $d$ os descontos. Fretes compartilhados
devem ser rateados ou registrados uma única vez para evitar duplicidade.

\begin{tabularx}{\textwidth}{@{}p{3.7cm}Y@{}}
\toprule
Campo do lançamento & Conteúdo esperado \\
\midrule
Identificador e categoria & Servidor, rede, backup ou serviço \\
Item e especificação & Modelo completo e quantidade \\
Fornecedor e datas & Cotação, compra, recebimento e pagamento \\
Preço e adicionais & Preço unitário, frete, encargos e descontos \\
Previsto e contratado & Estimativa e valor efetivamente acordado \\
Pago e saldo & Valores pagos, parcelas e valor ainda devido \\
Situação e evidência & Cotado, comprado, recebido, testado; comprovante e garantia \\
\bottomrule
\end{tabularx}

\textbf{Compras nesta versão:} o início da aquisição foi informado e a lista
de preços foi recebida, mas ainda faltam datas, fornecedores, situação dos
pedidos e valores pagos. As despesas realizadas serão inseridas quando esses
dados forem fornecidos. Em compras parceladas, o valor total contratado
integra o investimento; as parcelas descrevem o fluxo de caixa e não são somadas
novamente ao total.

\section{Consolidação do investimento}
\begin{align}
I_{\mathrm{base}} &= I_{\mathrm{PC}} + I_{\mathrm{infra}} + I_{\mathrm{servidor}}
+ I_{\mathrm{rede}} + I_{\mathrm{backup}},\\
I_{\mathrm{previsto}} &= I_{\mathrm{base}} + R,\\
S_{\mathrm{orcamento}} &= B - I_{\mathrm{previsto}},\\
\Delta I &= I_{\mathrm{realizado}} - I_{\mathrm{base}}.
\end{align}
$I_{\mathrm{PC}}$ representa a compra do computador principal e
$I_{\mathrm{infra}}$ os itens de rack e gabinetes.
$I_{\mathrm{servidor}}$ representa as novas despesas de adequação do servidor,
sem somar novamente o valor das peças antigas reaproveitadas ou os itens
já incluídos na infraestrutura física.
$I_{\mathrm{rede}}$ inclui o switch de \dinheiro{400,00}, o cabeamento e
os demais itens de rede que vierem a ser cotados; somente o switch tem preço
informado nesta versão.
$R$ representa a reserva para imprevistos e $B$ o limite informado,
nesta versão igual a \dinheiro{15.000,00}. A reserva integra esse limite,
não constitui autorização para ultrapassá-lo. Para comparar
previsto e realizado, os dois valores devem abranger o mesmo conjunto de itens.
Valores contratados, valores pagos e compras ainda previstas serão apresentados
separadamente. O saldo a pagar não é equivalente ao saldo do orçamento.

\section{Custos recorrentes e energia}
Os custos recorrentes poderão incluir energia, serviços contratados,
armazenamento externo e manutenção. O plano de internet existente só entra
como despesa incremental se houver aumento atribuível ao projeto; uma eventual
parcela do custo atual deverá ser apresentada como rateio separado.

O consumo mensal é estimado a partir da potência média medida na tomada:
\begin{align}
E_{\mathrm{mes}} &= \frac{P_{\mathrm{media}}}{1000}\,h_{\mathrm{dia}}\,n_{\mathrm{dias}},\\
C_{\mathrm{energia}} &= E_{\mathrm{mes}}\,T.
\end{align}
$P_{\mathrm{media}}$ está em watts, $h_{\mathrm{dia}}$ em horas por dia e $T$ é
a tarifa efetiva por kWh aplicável ao consumo incremental, com os componentes
variáveis pertinentes à conta. As potências nominais de 650 W da fonte existente
e de 850 W da fonte proposta não são medições do consumo do computador.
Equipamentos de rede adicionais também
devem ser considerados no limite de medição ou registrados separadamente.

\begin{tabularx}{\textwidth}{@{}YYY@{}}
\toprule
Potência hipotética & Consumo: 24 h/dia, 30 dias & Custo na tarifa $T$ \\
\midrule
30 W & 21,6 kWh & $21{,}6\,T$ reais \\
60 W & 43,2 kWh & $43{,}2\,T$ reais \\
100 W & 72,0 kWh & $72{,}0\,T$ reais \\
\bottomrule
\end{tabularx}
\textbf{Os três cenários são exemplos de cálculo, não medições ou previsões
do consumo deste servidor.} A tarifa e a potência real ainda não foram informadas.

\section{Custo total de propriedade}
Para um horizonte de $N$ meses, o cenário previsto pode ser expresso por:
\begin{equation}
\mathrm{TCO}(N)=I_{\mathrm{inicial}}+
\sum_{m=1}^{N}\left(C_{\mathrm{energia},m}+C_{\mathrm{servicos},m}
+C_{\mathrm{manutencao},m}\right).
\end{equation}
Reposições e expansões deverão ser incluídas uma única vez na categoria
correspondente. Poderão ser comparados cenários de 12 e 36 meses após obter
preços e medições. Não há estimativa de retorno financeiro nesta versão,
pois nenhum benefício monetário foi quantificado.

\chapter{Plano de execução, testes e riscos}
\section{Marcos do projeto}
\begin{tabularx}{\textwidth}{@{}p{0.8cm}p{4cm}YY@{}}
\toprule
ID & Marco & Dependência & Situação \\
\midrule
M1 & Finalizar a aplicação & Validação da versão de entrega & Em andamento \\
M2 & Confirmar inventário e orçamento & Modelos, valores e limite & Em andamento \\
M3 & Comprar e montar o computador & M2 e recebimento das peças & A documentar \\
M4 & Instalar e configurar o Proxmox VE & M3 e definição de discos e rede & Planejado \\
M5 & Criar a máquina virtual e preparar os serviços & M4 e escolha do sistema Linux & Planejado \\
M6 & Configurar rede e acesso & Levantamento e ambiente virtual disponível & Planejado \\
M7 & Implantar e validar a aplicação & M1, M5 e M6 & Planejado \\
M8 & Medir custos e iniciar operação & Backups e testes de aceitação & Planejado \\
\bottomrule
\end{tabularx}
Datas e duração serão definidas conforme a conclusão da aplicação, o orçamento
e a entrega dos componentes. O planejamento não atribui prazos ainda não acordados.

\section{Registro dos testes}
\begin{tabularx}{\textwidth}{@{}p{3.5cm}YY@{}}
\toprule
Área & Evidência a coletar & Resultado atual \\
\midrule
Hardware & Memória, discos, estabilidade e temperaturas & Sem medição registrada \\
Proxmox VE & Acesso administrativo, recursos, rede e reinício da máquina virtual & A executar \\
Aplicação & Login, contas, transações e isolamento & A validar na versão instalada \\
Serviços & Inicialização e recuperação após reinício & A executar \\
Rede & Acesso permitido, restrições e transferência & A executar \\
Backup & Integridade e restauração isolada & A executar no servidor \\
Energia & Potência média e período de medição & A executar \\
\bottomrule
\end{tabularx}
Cada resultado deverá informar ferramenta ou método, condições, data,
resultado esperado, resultado obtido e eventual correção aplicada.

\section{Riscos e decisões pendentes}
\begin{tabularx}{\textwidth}{@{}p{4cm}YY@{}}
\toprule
Situação & Efeito no projeto & Ação prevista \\
\midrule
CPU e placa-mãe incompatíveis na lista & Compra inadequada e atraso & Confirmar modelos antes da aquisição \\
DDR4 na lista Ryzen 9700X / B650 & Memória incompatível e alteração da estimativa & Corrigir para DDR5 e confirmar novo preço \\
Estado das peças desconhecido & Falhas ou reposições inesperadas & Inspeção e testes antes da adoção \\
Lista parcial e adequações pendentes & Custo final ainda indeterminado & Cotar cabeamento e demais itens dentro da margem de \dinheiro{1.756,00} \\
Backup só no próprio servidor & Perda conjunta de dados e cópias & Definir destino externo e testar restauração \\
Configuração de rede não levantada & Acesso ou desempenho inadequados & Mapear rede antes de comprar equipamentos \\
Procedimentos de deploy em evolução & Instalação incompleta & Conferir artefatos e validar no destino \\
\bottomrule
\end{tabularx}

\chapter{Diário e próximas atualizações}
\section{Registro inicial}
\begin{tabularx}{\textwidth}{@{}p{2.6cm}Y@{}}
\toprule
Data & Registro \\
\midrule
08/10/2026 & Criação da versão 0.1. Aplicação em finalização, início da aquisição
de componentes informado, inventário preliminar registrado e estrutura
financeira preparada. Orçamento global de \dinheiro{15.000,00} informado;
lista de compras a enviar posteriormente. Nenhuma montagem ou implantação
atestada nesta versão. \\
08/10/2026 & Versão 0.2. Nome atualizado para Projeto Finanças em todo o
documento. Proxmox VE incluído no planejamento, com instalação no computador
físico e máquina virtual Linux para Docker; configuração e testes pendentes. \\
08/10/2026 & Versão 0.3. Lista de oito componentes e valores registrada,
com subtotal de \dinheiro{11.100,00} e margem preliminar de \dinheiro{3.900,00}.
Memória DDR4 incompatível com AM5 registrada como pendência. Confirmada a
destinação da lista ao computador principal, com peças antigas para o servidor.
Arquivo principal definido como \path{relatorio/ProjetoFinancas.tex};
compras realizadas ainda a confirmar. \\
08/10/2026 & Versão 0.4. Incluídos rack de piso 19 polegadas, 12U, 470 mm,
por \dinheiro{944,00}, e dois gabinetes K-mex CR-S350, 4U, sem fonte,
considerando \dinheiro{400,00} por unidade. Total dos itens listados atualizado
para \dinheiro{12.844,00}, com margem preliminar de \dinheiro{2.156,00}.
Preço unitário dos gabinetes e encaixe físico a conferir. \\
08/10/2026 & Versão 0.5. Incluídos switch TP-Link TL-SG1016D por
\dinheiro{400,00} e cabeamento para dez máquinas, ainda sem metragem e preço.
Total dos itens com preço atualizado para \dinheiro{13.244,00}, com margem
preliminar de \dinheiro{1.756,00} antes do cabeamento e demais custos pendentes. \\
\bottomrule
\end{tabularx}

\section{Modelo para novas entradas}
% Copiar esta estrutura para cada evento real e preencher com fatos e evidências.
\begin{description}
\item[Data e etapa:] a preencher.
\item[Objetivo:] o que se pretendia realizar.
\item[Decisão e justificativa:] alternativa escolhida e motivo.
\item[Execução:] atividades efetivamente realizadas.
\item[Custo:] estimativa, valor contratado, valor pago e comprovante.
\item[Resultado e evidência:] teste, medição, fotografia ou registro técnico.
\item[Pendências:] próximo passo e dependências.
\end{description}

\section{Informações para a versão seguinte}
\begin{enumerate}
\item Nome do autor e eventual finalidade acadêmica ou institucional do relatório.
\item Conferência das peças antigas e sequência de transferência para o servidor.
\item Correção da memória para DDR5, preço atualizado e modelo completo da ASUS B650M.
\item Itens adicionais, fretes, distribuição da margem e gastos já realizados.
\item Conferência do preço unitário dos gabinetes, profundidade útil do rack e fixações.
\item Modelo correto da placa-mãe antiga e identificação completa das demais peças.
\item Versão final da aplicação e resultados de sua validação.
\item Equipamentos de rede existentes e local previsto para o servidor.
\item Percursos e metragens dos dez pontos, ligação ao roteador e cotação do cabeamento.
\item Versão do Proxmox VE, armazenamento e recursos previstos para a máquina virtual.
\end{enumerate}

\begin{thebibliography}{9}
\addcontentsline{toc}{chapter}{Referências}
\bibitem{repo} Projeto Finanças / FinanceHub.
\textit{Documentação local do repositório}: README da raiz, guia do backend,
arquitetura do frontend e guia de produção em home server.
Consultados em 8 out. 2026. Documentos de trabalho sujeitos a atualização.

\bibitem{intelcpu} Intel.
\textit{Intel Core i5 Desktop Processors: graphics and specifications}.
Disponível em
\url{https://www.intel.com/content/dam/support/us/en/documents/processors/core/Intel-Core-i5-HD-GFX.pdf}.
Acesso em: 8 out. 2026.

\bibitem{intelf} Intel.
\textit{Does My Computer with Intel Core Processor Number Ending with F Need a Discrete Graphics Card?}
Disponível em
\url{https://www.intel.com/content/www/us/en/support/articles/000058487/processors.html}.
Acesso em: 8 out. 2026.

\bibitem{amdb450} AMD.
\textit{AMD Socket AM4 Chipsets: Specs and CPU Compatibility}.
Disponível em \url{https://www.amd.com/en/products/processors/chipsets/am4.html}.
Acesso em: 8 out. 2026.
\bibitem{proxmox} Proxmox Server Solutions.
\textit{Proxmox VE Administration Guide}.
Disponível em \url{https://pve.proxmox.com/pve-docs/pve-admin-guide.html}.
Acesso em: 8 out. 2026.
\bibitem{ryzen9700} AMD.
\textit{AMD Ryzen 7 9700X: especificações}.
Disponível em
\url{https://www.amd.com/en/products/processors/desktops/ryzen/9000-series/amd-ryzen-7-9700X.html}.
Acesso em: 8 out. 2026.

\bibitem{asusb650} ASUS.
\textit{PRIME B650M-A: especificações técnicas}.
Referência da plataforma B650; não identifica o modelo exato da lista.
Disponível em
\url{https://www.asus.com/motherboards-components/motherboards/prime/prime-b650m-a/techspec/}.
Acesso em: 8 out. 2026.
\bibitem{tplinkswitch} TP-Link Brasil.
\textit{TL-SG1016D: switch Gigabit de mesa ou rack de 16 portas}.
Disponível em
\url{https://www.tp-link.com/br/business-networking/soho-switch-unmanaged/tl-sg1016d/}.
Acesso em: 8 out. 2026.
\end{thebibliography}
\end{document}
